% concatenated from The_Capillary_Lp_Dual_Minkowski_Problem_Sections_4_5_Revised_20260907.tex

\documentclass[11pt]{amsart}
%\usepackage{showkeys}


\textwidth 16.00cm
\textheight 20cm
\topmargin 0.0cm
\oddsidemargin 0.0cm

\evensidemargin 0.0cm
\parskip 0.0cm


%\textwidth=5.4in \textheight=8in




\usepackage{amssymb}
\usepackage{mathrsfs}
\usepackage{amsfonts}
\usepackage{latexsym,amsmath,amsthm,amssymb,amsfonts}
\usepackage[usenames]{color}
\usepackage{xspace,colortbl}
\usepackage{graphicx}
\usepackage{tipa}




\newcommand{\be}{\begin{equation}}
\newcommand{\ee}{\end{equation}}
\newcommand{\beq}{\begin{eqnarray}}
\newcommand{\eeq}{\end{eqnarray}}
\newcommand{\cM}{\mathcal{M}}
\newcommand{\ud}{\mathrm{d}}
\newcommand{\re}{\mathrm{e}}
\newcommand{\wh}{\tilde{h}}
\newcommand{\rt}{\mathrm{tr}}
\newcommand{\pr}{\partial}
\newcommand{\fr}{\frac}
\newcommand{\uS}{\mathbb{S}^{n}}

\def\A{{\mathfrak A}}
\def\B{{\mathfrak B}}
\def\C{{\mathfrak C}}
\def\W{{\mathfrak W}}
\def\D{{\mathfrak D}}
\def\E{{\mathfrak E}}
\def\F{{\mathfrak F}}
\def\f{{\mathfrak f}}
\def\G{{\mathfrak G}}
\def\g{\mathfrak g}
\def\H{{\Bbb H}}
\def\h{{\mathfrak h}}
\def\M{{\mathfrak M}}
\def\P{{\mathfrak P}}
\def\R{{\mathfrak R}}
\def\U{{\mathfrak U}}
\def\fS{{\mathfrak S}}
\def\W{{\mathfrak W}}
\def\V{{\mathfrak V}}
%\def\{{\mathfrak }}
%\def\{{\mathfrak }}
%\def\{{\mathfrak }}
%\def\{{\mathfrak }}
%\def\C{{\mathfrak C}}

\def\cA{{\mathcal A}}
\def\cC{{\mathcal C}}
\def\cB{{\mathcal B}}

\def\cD{\mathcal D}

\newtheorem{prop}{Proposition}[section]
\newtheorem{theo}[prop]{Theorem}
\newtheorem{lemm}[prop]{Lemma}
\newtheorem{coro}[prop]{Corollary}
\newtheorem{rema}[prop]{Remark}
\newtheorem{remark}[prop]{Remark}
\newtheorem{exam}[prop]{Example}
\newtheorem{defi}[prop]{Definition}
\newtheorem{conj}[prop]{Conjecture}
\newtheorem{pro}[prop]{Problem}
\newtheorem{RK}{Remark}
\def\begeq{\begin{equation}}
\def\endeq{\end{equation}}
\def\p{\partial}
\def\cF{\mathcal F}
\def\lf{\left}
\def\ri{\right}
\def\e{\epsilon}
\def\ol{\overline}
\def\R{\mathbb R}
\def\S{\mathbb S}
\def\wt{\widetilde}
\def\tr{{\rm tr}}
\def\hbk{\mathbb{H}^{ 3}_{-\kappa^2 }}
\def\hb{\mathbb{H}^3_{-1}}
\def\ii{\sqrt{-1}}
\def\la{\langle}
\def\ra{\rangle}
\def\Minkowski{\mathbb{R}^{3,1}}

\def\d{\delta}
\def\a{\alpha}
\def\b{\beta}
\def\s{\sigma}

\def \ds{\displaystyle}
\def \vs{\vspace*{0.1cm}}

%\def\S{\mathbb S}
\def\T{\mathbb T}

\def\odot{\setbox0=\hbox{$\bigcirc$}\relax \mathbin {\hbox
to0pt{\raise.5pt\hbox to\wd0{\hfil $\wedge$\hfil}\hss}\box0 }}

\numberwithin{equation} {section}


\numberwithin{equation}{section}
\textheight=8.8in
\textwidth=6.28in
\topmargin=0mm
\oddsidemargin=0mm
\evensidemargin=0mm


\def\endproof{$\hfill\Box$\\}
\def\s{\,\,\,\,}
\def\mv{2.0ex}
\def\Re{\mbox{Re}}
\newtheorem{theorem}{\bf Theorem}[section]
\newtheorem{proposition}[theorem]{\bf Proposition}
\newtheorem{definition}[theorem]{\bf Definition}
%\newtheorem{remark}[theorem]{\bf Remark}
\newtheorem{lemma}[theorem]{\bf Lemma}
\newtheorem{example}[theorem]{\bf Example}
\newtheorem{corollary}[theorem]{\bf Corollary}
\def\dint{\displaystyle{\int}}
\allowdisplaybreaks

\begin{document}
\title[The capillary $L_p$ dual Minkowski problem]
 {The capillary $L_p$ dual Minkowski problem for $p>q$ in higher dimensions}

\author{ Ya Gao$^{\dagger,\ast}$}

\address{
 $^{\dagger}$School of Mathematics and Statistics, South-Central Minzu University, Wuhan 430074, China. }

\email{Echo-gaoya@outlook.com}

\thanks{$\ast$ Corresponding author}

\date{}
%\maketitle
\begin{abstract}
We prove existence and uniqueness for the capillary $L_p$ dual Minkowski problem for $p>q$ and contact angle $\theta\in (0,\frac{\pi}{2})$ in $\R^{n+1}$, $n\geq 3$. We reduce it to a Monge-Amp\`ere type equation with a Robin boundary condition on the unit spherical cap, by building a new auxiliary function to obtain the $C^2$ estimate for $p>q$ in $\R^{n+1}$ ($n\geq 3$), we then prove that there exists a unique smooth solution that solves this problem provided $\theta\in (0, \frac{\pi}{2})$.
\end{abstract}

\maketitle {\it \small{{\bf Keywords}: $L_p$ dual Minkowski problem, Capillary hypersurfaces, Monge-Amp\`ere equation, Robin boundary condition.}

{{\bf MSC 2020}: Primary 53C65, 35J66. Secondary 53C42, 53C45, 35J60.}}





\section{Introduction}
The classical Minkowski problem is a fundamental problem in the Brunn--Minkowski theory of convex geometry. It asks one to determine a convex body whose surface area measure corresponds to a prescribed spherical Borel measure. This problem has been completely solved under a necessary and sufficient condition by Minkowski \cite{Min97, Min03}, Alexandrov \cite{Ale38, Ale39, Ale56}, Lewy \cite{Lew38}, Nirenberg \cite{Nir53}, Pogorelov \cite{Pog52}, Cheng-Yau \cite{CY76}, and others. 

The classical Brunn--Minkowski theory has been extended by the $L_p$ Brunn--Minkowski theory and the dual Brunn--Minkowski theory, and these theories are now central in convex geometry. The $L_p$-Minkowski problem is a fundamental problem in the $L_p$ Brunn--Minkowski theory and greatly generalizes the classical Minkowski problem. It was introduced by Lutwak \cite{Lut.JDG.38-1993.131, Lut.Adv.118-1996.244} and has been extensively studied since then; see e.g.
\cite{%BHZ.IMRNI.2016.1807,
  BLYZ.JAMS.26-2013.831,
  CLZ.TAMS.371-2019.2623,
  Zhu.Adv.262-2014.909}
for the logarithmic Minkowski problem,
\cite{JLW.JFA.274-2018.826,
  JLZ.CVPDE.55-2016.41,
  Lu.SCM.61-2018.511,
  Lu.JDE.266-2019.4394,
  LW.JDE.254-2013.983,
  Zhu.JDG.101-2015.159}
for the centroaffine Minkowski problem, and
\cite{CW.Adv.205-2006.33,
  HLYZ.DCG.33-2005.699,
  % LYZ.JDG.56-2000.111,
  LYZ.TAMS.356-2004.4359,
  Sta.Adv.167-2002.160}
for other cases of the $L_p$-Minkowski problem. On the other hand, the dual Minkowski problem was first proposed by Huang, Lutwak, Yang and Zhang
in their recent groundbreaking work \cite{HLYZ.Acta.216-2016.325} and then
followed by
\cite{BHP.JDG.109-2018.411,
  %CL.Adv.333-2018.87,
  HP.Adv.323-2018.114,
  HJ.JFA.277-2019.2209,
  %JW.JDE.263-2017.3230,
  LSW.JEMSJ.22-2020.893,
  Zha.CVPDE.56-2017.18,
  Zha.JDG.110-2018.543}.

  Recently, Lutwak, Yang and Zhang introduced the $L_p$ dual Minkowski problem in \cite{LYZ18}, which unifies the classical Minkowski problem, the $L_p$-Minkowski problem and the dual Minkowski problem. For the general $L_p$ dual Minkowski problem, much progress has already been made
\cite{BF.JDE.266-2019.7980,
  CHZ.MA.373-2019.953,
  CCL,
  % CL,
  HLYZ.JDG.110-2018.1,
  HZ.Adv.332-2018.57,
  LLL}. 
  In fact, for a convex body $K\subset \mathbb{R}^{n+1}$ containing the origin in its interior, the $L_p$ dual curvature measure $d\widetilde{C}_{p,q}$\footnote{Sometimes simply referred to as \emph{the $(p,q)$-th dual curvature measure}.} is defined by
  \begin{equation*}
      d\widetilde{C}_{p,q} =\frac{1}{n+1}h_K^{1-p}\left(h_K^2 + |\nabla h_K|^2\right)^{\frac{q-n-1}{2}}\det(\nabla^2 h_K+h_K\delta_{ij})da,
  \end{equation*}
where $h_K$ is the support function of $K$, $\nabla h_K$ and $(h_K)_{ij}$ are the gradient and the Hessian of $h_K$ on the unit sphere $\uS$ with respect to an orthonormal basis respectively, 
$da$ is the standard spherical area measure.\\
{\bf The $L_p$ dual Minkowski problem.} \emph{Given a finite nonzero Borel measure $m$ on $\uS$ and real numbers $p,q$, does there exist a convex hypersurface (as a boundary of a convex body) such that its induced area measure $d\tilde{C}_{p,q}$ equals $m$?}

When the given measure $m$ has a density $f$, the $L_p$ dual Minkowski problem becomes the existence problem for the following Monge-Amp\`ere equation on $\uS$:
$$\det(h_{ij}+ h\delta_{ij} ) = fh^{p-1}(h^2 + |\nabla h|^2)^{\frac{n+1-q}{2}},$$
where $f$ is a given positive smooth function on $\uS$, $\delta_{ij}$ is the Kronecker delta, $\nabla h$ and $(h_{ij})$ are the gradient and the Hessian of $h$ on $\uS$ with respect to an orthonormal basis respectively.

From the above Monge-Amp\`ere equation, when $q=n+1$ the $L_p$ dual Minkowski problem becomes the $L_p$ Minkowski problem, which includes the classical Minkowski problem; when $p=0$, the $L_0$ dual Minkowski problem becomes the dual Minkowski problem; if, in addition, $q=0$, the dual Minkowski problem becomes the Aleksandrov problem.

As the classical Minkowski problem and related problems have developed, analogous boundary problems have naturally emerged. Busemann \cite{Busemann59} considered Minkowski and related problems for convex surfaces with boundaries early on. Oliker \cite{Oli82} introduced and studied a boundary problem equivalent to solving the Monge-Amp\`ere equation of the classical Minkowski problem with homogeneous Dirichlet boundary condition. Other Dirichlet boundary value problems have also been studied in \cite{CW95, Pog78, Sch18, Sch03}. For the analytic theory of Neumann problems for equations of Monge-Amp\`ere type, see \cite{LTU86}. Recently, several works have introduced Robin or Neumann boundary value conditions for convex capillary hypersurfaces.

\subsection{Setup and the problem }

Let $\{E_i\}_{i=1}^{n+1}$ be the standard orthonormal basis of $\R^{n+1}$, $\R^{n+1}_{+}=\{x\in\R^{n+1}| x\cdot E_{n+1}>0\}$ be the upper Euclidean half-space, $\Sigma\subset\overline{\R^{n+1}_{+}}$ be a properly embedded, smooth compact hypersurface with boundary such that
$$int(\Sigma)\subset\R^{n+1}_{+}\qquad and \qquad\partial\Sigma\subset\partial\R^{n+1}_{+}.$$
We call $\Sigma\subset\overline{\R^{n+1}_{+}}$ a {\emph{capillary hypersurface}} if $\Sigma$ intersects $\partial\R^{n+1}_{+}$ at a constant contact angle $\theta\in (0,\pi)$. Let $\nu$ be the unit outward normal, i.e., the Gauss map of $\Sigma$ with respect to the domain $\hat{\Sigma}$, where $\hat{\Sigma}$ is a bounded closed region in $\overline{\R^{n+1}_{+}}$ enclosed by the convex capillary hypersurface $\Sigma$ and $\partial\R^{n+1}_{+}$, and denote $\hat{\partial\Sigma}:=\partial\hat{\Sigma}\backslash\Sigma\subset\partial\R^{n+1}_{+}$. The contact angle $\theta$ is defined by 
$$\cos(\pi-\theta) = \langle \nu , e\rangle,$$
where $e:= -E_{n+1}$, hence $e$ is the unit outward normal of $\partial\R^{n+1}_{+}$. If $\Sigma$ is convex, then the Gauss image $\nu(\Sigma)$ of $\Sigma$ lies in the spherical cap
$$\S^n_{\theta}:=\left\{x\in\uS | \langle x, E_{n+1}\rangle \geq \cos\theta\right\}.$$
Instead of the usual Gauss map $\nu$, we give another map
$$\tilde{\nu}:= T \circ \nu: \Sigma\to C_{\theta},$$
where $C_{\theta}$ is a spherical cap defined by
$$C_{\theta}:= \left\{\xi\in\overline{\R^{n+1}_{+}} |~ |\xi-\cos\theta\cdot e|=1\right\},$$
which is also a capillary hypersurface, and $T: \S^n_{\theta}\to C_{\theta}$ is the translation in the vertical direction defined by $T(z)=z+\cos\theta\cdot e$. The diffeomorphism $\tilde{\nu}$ is called the {\emph{capillary Gauss map}} of $\Sigma$, and thus we can reparametrize $\Sigma$ using its inverse on $C_{\theta}$ (we can refer to \cite[Fig. 1]{MWW25a}). Hence, we also view the usual support function $h$ of $\Sigma$ as a function defined on $C_{\theta}$. 
Denote $\mathcal{K}_{\theta}$ as the set of all capillary convex bodies in $\overline{\R^{n+1}_{+}}$, and $\mathcal{K}^{\circ}_{\theta}$ as the family of capillary convex bodies with the origin as an interior point of the flat part of their boundary.

Recently, Mei, Wang and Weng \cite{MWW25a} introduced a capillary Minkowski problem, which asks for the existence of a strictly convex capillary hypersurface $\Sigma\subset\overline{\R^{n+1}_{+}}$ with prescribed Gauss-Kronecker curvature on a spherical cap $C_{\theta}$. Subsequently, they considered a capillary $L_p$-Minkowski problem for $p\geq 1$ in \cite{MWW25b}. Hu, Ivaki and Scheuer \cite{HIS25} considered the capillary Christoffel-Minkowski problem, while Hu and Ivaki \cite{HI25} solved the even capillary $L_p$-Minkowski problem for the range $-n< p < 1$ and $\theta\in (0, \frac{\pi}{2})$ using an iterative scheme. Wang and Zhu \cite{WZ25} introduced the more general capillary Orlicz-Minkowski problem.

In addition to this, we \cite{GY26v2} consider the $L_p$ dual Minkowski problem for $p>q$ and $q\leq 1$ in $\R^{n+1}$. Due to the selection of the auxiliary function for the $C^2$ estimate, we are unable to eliminate the condition of $q\leq 1$. Recently, Hu and Yang \cite{HY26} refined $C^2$ argument uses the auxiliary function $\mathfrak{Q}=\log\sigma_1+\mathfrak{c}_1 h+\mathfrak{c}_2 |\nabla h|^2$ and prove the capillary $L_p$ dual Minkowski problem in the case $p>q$ without requiring $q\leq 1$ in $\R^3$, but they observe that this test function does not yield the required $C^2$ estimate in dimensions $n\geq 3$. 

A complementary approach uses curvature flow as a tool. From this point of view, Mei, Wang and Weng \cite{MWW25c} studied a Gauss curvature type flow for capillary hypersurfaces, namely the capillary Gauss curvature flow. Hu, Hu and Ivaki \cite{HHI25} obtained the long-time existence and asymptotic behavior of a class of anisotropic capillary Gauss curvature flows, and as applications their flow gives smooth solutions to the capillary even $L_p$ Minkowski problem in the Euclidean half-space and to the capillary $L_p$ Minkowski problem. 

In this paper, we want to study the {\emph{capillary $L_p$ dual Minkowski problem}}, i.e., study the capillary $(p,q)$-th dual curvature measure $d\tilde{C}^c_{p,q}$ for a convex capillary body in $\overline{\mathbb{R}_+^{n+1}}$, which is defined by
$$d\tilde{C}^c_{p,q}:= \frac{1}{n+1}lh^{1-p}\left(h^2 + |\nabla h|^2\right)^{\frac{q-n-1}{2}}\det(\nabla^2 h+h\sigma)d\sigma,$$
where $\nabla h$ and $\nabla^2 h$ are the gradient and the Hessian of $h$ on $C_{\theta}$ with respect to the standard spherical metric $\sigma$ on $C_{\theta}$ respectively, and $l:=\sin^2\theta + \cos\theta\langle\xi,e\rangle$, $\xi\in C_{\theta}$.\\ 
\\
{\bf Capillary $L_p$ dual Minkowski problem.} \emph{Given a positive smooth function $f$ on $C_{\theta}$, does there exist a capillary convex body $\Sigma\in \mathcal{K}^{\circ}_{\theta}$ such that its capillary $(p,q)$-th dual curvature measure $d\tilde{C}^c_{p,q}$ equals $f\,l\,d\sigma$?}

By applying a similar argument to that in \cite[Proposition 2.4]{MWW25a}, the capillary $L_p$ dual Minkowski problem is equivalent to solving the following Neumann boundary value problem for a Monge-Amp\`ere type equation:
\begin{equation}\label{Eq}
\left\{
\begin{aligned}
&\det\left(\nabla^2 h + h\sigma\right)=fh^{p-1}\left(h^2 + |\nabla h|^2\right)^{\frac{n+1-q}{2}}, ~~&&in~~C_{\theta}\\
&\nabla_{\mu} h=\cot\theta\cdot h,  ~~&& on~~\partial C_{\theta}
\end{aligned}
\right.
\end{equation}
where $\mu$ is the unit outward normal of $\partial C_{\theta}\subset C_{\theta}$. Here, if 
$A:=\nabla^2 h + h\sigma >0 ~~in~C_{\theta}$
is a solution of Eq. \eqref{Eq}, then we call it a convex solution.

\subsection{Main results}

We consider the capillary $L_p$ dual Minkowski problem for $p>q$ (without requiring $q\leq 1$) in $\R^{n+1}$, $n\geq 3$, as formulated in Eq. \eqref{Eq}. The case $q=n+1$ is the capillary $L_p$-Minkowski problem; for $p>q=n+1$ it has already been solved in \cite{MWW25b}; for $p>q, q\leq 1$ and $p>q, n=2$ have been solved by \cite{GY26v2} and \cite{HY26} respectively. The main theorem of this paper extends the previous conclusion to $L_p$ dual Minkowski problem for $p>q$ and $n\geq 3$. It stated below.

\begin{theorem}\label{main1.1}
Let $n\geq 3$, $p>q$ and $\theta\in (0, \frac{\pi}{2})$. For any positive smooth function $f$ defined on $C_{\theta}$, there exists a unique positive smooth strictly convex solution $h$ of Eq. \eqref{Eq}. 
\end{theorem}

To prove the main theorem, we establish a priori estimates for solutions to Eq. \eqref{Eq} up to the second derivative, and then apply the continuity method to obtain existence. First, we give uniform positive lower and upper bounds for solutions to prevent Eq. \eqref{Eq} from degenerating.

When $p>q$, we use the capillary support function $u(\xi)$, which satisfies a Monge-Amp\`ere type equation with vanish Neumann boundary condition. The maximum principle gives positive lower and upper bounds for the solutions; see Lemma \ref{C0 estimate}. Then we give a rough $C^1$ estimate in Lemma \ref{c1-1}, based only on convexity and independent of the specific equation. Namely, the gradient of the solution $h$ is bounded by the contact angle $\theta$ and the upper bound of $h$ as follows:
$$\max_{C_{\theta}}|\nabla h|\leq (1+\cot^2\theta)^{1/2}\Vert h\Vert_{C^0 (C_{\theta})},$$
so we have the upper bound of $|\nabla h|$. 

For the $C^2$ estimate, the existing text functions are unable to handle the cases in high dimensions for $p>q$ without $q\leq 1$. Inspired by the auxiliary function $\mathfrak{Q}=\log\sigma_1 +\mathfrak{c}_1h +\mathfrak{c}_2 | \nabla h|^2$ in Hu-Yang's second version on arXiv \cite{HY26}, we consider a new text function as follows
$$Q:=\log\sigma_1(A)+a h+b|\nabla h|^2+\eta\log\sigma_{n-1}(A),$$
where the constant $b>0$, $a<0$, and in the final proof we take $\eta=1$. %%The identity
%%$$\sigma_{n-1}(A)=\det(A)\rt(A^{-1})$$
%%shows why the extra term is useful: under the equation, it detects the higher-dimensional anisotropic collapse in which the determinant is fixed while some tangential eigenvalues become small. The proof of the $C^2$ estimate first records the differential identities for $\log\sigma_{n-1}(A)$, then proves the interior and boundary maximum estimates for this extended $P$, and finally converts the bound for $P$ into a bound for $\sigma_1(A)$.

\begin{remark}
    %%{\bf (a)} The case $q=n+1$ is included in Theorem \ref{main1.1}, but this is the capillary $L_p$-Minkowski problem and is already covered by \cite{MWW25b} for $p>n+1$;\\
    {\bf (a)}
    The range $\theta\in (0,\frac{\pi}{2})$ is used in the global $C^2$ estimate, because it gives the strict convexity of $\partial C_{\theta}\subset C_{\theta}$. This restriction on the range is not necessary in establishing the $C^0$ and $C^1$ estimates for Eq. \eqref{Eq};\\ 
    {\bf (b)}
    When $n=2$, the auxiliary function we have chosen reduces to Hu-Yang's \cite{HY26} text function. There, in this paper, we only prove the case where $n\geq 3$;\\
    {\bf (c)}
In fact, after completing this article, we know that Hu-Yang also has reselected the auxiliary function for $C^2$ estimate (different with text function $Q$ in this paper), it can also solve the capillary $L_p$ dual Minkowski problem for $p>q$ and $n\geq 3$, without relying on the conditions that $q\leq 1$ and $n=2$.

   %% {\bf (b)}
    %%The added $\log\sigma_{n-1}(A)$ term is the mechanism that replaces the former restriction $q\leq 1$ in the interior trace test.
\end{remark}


This paper is organized as follows. In Section 2, we recall the preliminaries needed below. In Section 3, we prove the $C^0$ and gradient estimate for solutions of Eq. \eqref{Eq}. In Section 4, we first present some notations, formulas and Lemmas that are necessary for the curvature estimate, and then prove the global $C^2$ estimate. In Section 5, we prove uniqueness and use the continuity method to establish Theorem \ref{main1.1}.



\section{Preliminaries } \label{se2}
This section is divided into two subsections. In the first subsection, we introduce the notion of capillary $L_p$ dual curvature measure, i.e., the capillary $(p,q)$-th dual curvature measure. 
In the second subsection, we provide some basic properties of capillary convex hypersurfaces.  

\subsection{The capillary $L_p$ dual Minkowski problem}

In this subsection, we recall the $L_p$ dual curvature measures from \cite{LYZ18} and describe their capillary counterpart.

Let $\Sigma$ be a capillary hypersurface in $\overline{\mathbb{R}^{n+1}_+}$. Its support function $h$ is defined by
$$h (y) = \max \{ x\cdot y : x\in \Sigma\}, \quad y\in\mathbb{R}^{n+1}.$$

Suppose that $\Sigma\in \mathcal{K}_{\theta}^{\circ}$. The radial function $\rho$ and radial map $r$ are defined by
\begin{equation*}
    \rho (x) = \max\{\lambda: \lambda x\in \Sigma\}, \qquad x\in \mathbb{R}^{n+1}\backslash\{0\}
\end{equation*}
and
\begin{equation*}
    r: \uS_+\to \Sigma, \qquad r(\gamma)=\rho(\gamma)\gamma \in\Sigma
\end{equation*}
respectively, where $\uS_+$ is the upper sphere. For $\gamma\in\uS_+$, define the \emph{radial Gauss map $\alpha:\uS_+\to \uS_{\theta}$ of $\gamma$} by 
\begin{equation*}
    \alpha(\gamma) = \nu(\rho (\gamma)\gamma)\in \uS_{\theta}
\end{equation*}
and $\alpha^{\ast}: \uS_{\theta} \to \uS_+$ is its reverse map. Similar to the definition of capillary Gauss map, instead of the usual radial Gauss map $\alpha$, we give another map 
\begin{equation*}
    \widetilde{\alpha}:= T\circ \alpha \circ r^{-1}  : \Sigma \to C_{\theta},
\end{equation*}
so its reverse map is $\widetilde{\alpha^{\ast}}: C_{\theta} \to \Sigma$. The map $\widetilde{\alpha}$ is called the \emph{capillary radial Gauss map} of $\Sigma$, and we therefore view the usual radial function $\rho$ of $\Sigma$ as a function on $C_{\theta}$.

For $K,L\in\mathcal{K}^{\circ}_{\theta}$ and $\lambda, \tau\geq 0$, the \emph{Minkowski combination} $\lambda K+ \tau L$ is defined by $\lambda K+ \tau L :=\{\lambda x+ \tau y: x\in K, y\in L\}$ and 
$$h_{\lambda K+ \tau L}:= \lambda h_{K} + \tau h_{L},$$
where $h_K$ and $h_L$ denote the support functions of the convex body $K$ and $L$ respectively.

The \emph{$L_p$ combination}, an extension of the Minkowski combinations studied by Firey in the early 1960s, is defined by
\begin{equation*}
    h^p_{\lambda K+_{p}\tau L }:= \lambda h_{K}^p + \tau h_{L}^p
\end{equation*}
for each $p\geq 1$, $K, L\in\mathcal{K}^{\circ}_{\theta}$ and $\lambda, \tau \geq 0$. 


The $(p,q)$-th dual curvature measure was introduced by Lutwak, Yang and Zhang in \cite{LYZ18}. Based on the definition in \cite{LYZ18} and the above introduction, we can define the capillary $(p,q)$-th dual curvature measure of $\Sigma\in \mathcal{K}_{\theta}^{\circ}$ as 
$$ d\widetilde{C}_{p,q}^c =\frac{1}{n+1}lh_{\Sigma}^{1-p}\left(h_{\Sigma}^2 + |\nabla h_{\Sigma}|^2\right)^{\frac{q-n-1}{2}}\det(\nabla^2 h_{\Sigma}+h_{\Sigma}\sigma)d\sigma.$$




The preceding definition gives a more general area measure for a capillary convex body $\Sigma\in \mathcal{K}^{\circ}_{\theta}$. If $q=n+1$, it coincides with the capillary $L_p$-surface area measure, and our result reduces to the case $p>n+1$ treated in \cite{MWW25b}.

\subsection{Basic properties of capillary convex hypersurfaces}
We first recall some basic properties of convex hypersurfaces in $\mathbb{R}^{n+1}$. Let $\Sigma\subset\overline{\mathbb{R}^{n+1}_{+}}$ be a smooth, properly embedded, strictly convex capillary hypersurface as above. We parametrize $\Sigma$ by the inverse capillary Gauss map $X: C_{\theta}\to \Sigma$ given by
$$X(\xi) = \tilde{\nu}^{-1}(\xi) = \nu^{-1}\circ T^{-1}(\xi) = \nu^{-1}(\xi-\cos\theta e), \quad \xi\in C_{\theta}.$$
The usual support function $h$ of $\Sigma$ is given by
$$h(X) : =\langle X, \nu(X)\rangle.$$
The support function $h$ is equivalently defined by the unique decomposition
$$X = h(X)\nu (X) + s, \qquad s\in T_{X}\Sigma.$$
Now, we define the capillary support function $u$ by the following unique decomposition
\begin{equation}\label{e1}
     X = u(X)\tilde{\nu}(X) + s', \qquad s'\in T_{X}\Sigma.
\end{equation}
It's easy to see that
$$h(X) = u(X)\langle\nu(X), \tilde{\nu}(X)\rangle = u(X)(1+\cos\theta\langle\nu,e\rangle)$$
since $\tilde{\nu} := T\circ \nu = \nu + \cos\theta e$.
Now, let 
$$h(\xi):= \langle X(\xi),\nu(X(\xi))\rangle = \langle X(\xi), T^{-1}(\xi)\rangle = \langle\tilde{\nu}^{-1}(\xi), \xi-\cos\theta e\rangle,$$
referring to \cite[Proposition 2.4]{MWW25a}, it satisfies 
$$\nabla_{\mu}h = \cot\theta h, \qquad on~ \partial C_{\theta},$$
and 
$$u(\xi) = \frac{h(\xi)}{|\xi|^2 - \cos\theta\langle \xi,e\rangle} = \frac{h(\xi)}{\sin^2\theta + \cos\theta\langle \xi,e\rangle} = \frac{h(\xi)}{l(\xi)},$$
due to $|\xi - \cos\theta e|^2 =1$.

Now we compute the geometric quantities of $X$ in terms of $h$ and $u$. Let $\{e_i\}_{i=1}^n$ be a local orthonormal frame on $C_{\theta}$ such that along the $\partial C_{\theta}$, $e_n = \mu$ is the unit outward normal of $\partial C_{\theta}\subset C_{\theta}$. Together with $T^{-1}(\xi) = \xi -\cos\theta e$, which is the usual normal of $C_{\theta}$ at $\xi$, it builds on a local orthonormal frame in $\R^{n+1}_{+}$. Hence by the definition of $h$, for $X(\xi)\in\Sigma$,
$$X = \sum_{i=1}^n\langle X, e_i\rangle e_i + \langle X, T^{-1}(\xi)\rangle T^{-1}(\xi) = \sum_{i=1}^n\langle X, e_i\rangle e_i + h(\xi)T^{-1}(\xi).$$
Since $\nabla$ is the standard connection in $C_{\theta}$ with respect to the standard spherical metric $\sigma$, it's easy to see
$$\sum_{i=1}^n \langle X, e_i\rangle e_i = \nabla h(\xi).$$
Hence
$$X(\xi) = \nabla h(\xi) + h(\xi)T^{-1}(\xi).$$
By the direct computation
$$\nabla_{e_j}X = (\nabla_{ij}h + h\delta_{ij})e_i,$$
where $\delta_{ij}$ is the coefficients of the standard metric $\sigma$ on $\uS$. Hence, the second fundamental form $A_{ij}$ of $\Sigma$ is given by
$$A_{ij} = \langle \nabla_{e_j}X, \nabla_{e_i}(\xi-\cos\theta e)\rangle = \langle \nabla_{e_j}X, e_i\rangle = \nabla_{ij}h + h\delta_{ij},$$
and the induced metric $g_{ij}$ of $\Sigma\subset\overline{\R^{n+1}_{+}}$ can be derived by Weingarten's formula
$$g_{ij} = \langle \nabla_{e_i}X, \nabla_{e_j}X\rangle = A_{ik}A_{jl}\delta^{kl}.$$
The Gauss-Kronecker curvature of $\Sigma$ at $X$ is 
$$K_{G}(X(\xi)) = \det (g^{ik}A_{kj}) = \det (\nabla ^2 h(\xi)+ h(\xi)\sigma)^{-1}.$$

Both functions and formula will play a crucial role in our paper. 
Next, we give the equation of $h$ and $u$ on $\partial C_{\theta}$, which will be used in subsequent proof process.

Along the boundary $\partial C_{\theta}$, we choose an orthonormal frame $\{e_{i}\}_{i=1}^n$ with $e_n = \mu$ and $\mu$ is the unit outer normal of $\partial C_{\theta}$. Then 

\begin{proposition}\label{pro}
    The support function $h$ and the capillary support funtion $u=l^{-1}h$ satisfy the following boundary conditions on $\partial C_{\theta}$:
    \begin{equation}\label{pro1}
        h_{kn}=0,
    \end{equation}
    which is equivalent to 
    \begin{equation}\label{pro2}
        u_{kn} = -\cot\theta u_{k},
    \end{equation}
    where $k=1,2,\cdots, n-1$.
\end{proposition}

\begin{proof}
    First, taking tangential derivative of the equation $h_n = \cot\theta h$ along $\partial C_{\theta}$, we have $\nabla_{e_k}h_n = \cot\theta h_k$. It follows that 
    \begin{equation*}
\begin{split}
    h_{kn} &= \nabla^2 h(e_k,e_n) = \langle \nabla_{e_k}(\nabla h), e_n\rangle  \\ 
    &= 
    \nabla_{e_k}(\langle \nabla h, e_n\rangle) - \langle\nabla h,\nabla _{e_k}e_n\rangle \\ 
    &= \nabla_{e_k}h_n -\langle \nabla h, \cot\theta e_k\rangle \\ 
    &= 
    \cot\theta h_k - \cot\theta h_k = 0,
\end{split}
    \end{equation*}
    where we have used $\nabla_{e_k}e_n = \cot\theta e_k$. Then the equation \eqref{pro1} follows. Using the same computation and the fact $u_n =0$ in \eqref{a2}, we also have \eqref{pro2}. We complete the proof.
\end{proof}



\section{A priori estimates} \label{se3}
In this section, we give a priori estimates (including $C^0$ and $C^1$ estimates) for solutions to Eq. \eqref{Eq}.

\begin{lemma}[{\bf $C^0$ estimate}]\label{C0 estimate}
    Let $p>q$ and $\theta\in(0, \pi)$. Suppose $h$ is a positive smooth strictly convex solution to Eq. \eqref{Eq}, that is, $A=\nabla^2h+h\sigma>0$ on $C_{\theta}$. Then we have 
    \begin{equation}\label{a1}
       \min\{ 1-\cos\theta, \sin^2\theta\}^{p-q}\cdot C_2
        \leq h^{p-q}\leq
       \max\{ 1-\cos\theta, \sin^2\theta\}^{p-q}\cdot C_1 ,
    \end{equation}
    where the positive constant $C_1$ and $C_2$ depend on $n$, $p$, $q$, $\theta$, $\min_{C_{\theta}}f$ and $\Vert f\Vert_{C^0}(C_{\theta})$.
    In particular, $h$ has positive lower and upper bounds depending only on
    $n,p,q,\theta$, $\min_{C_{\theta}}f$ and $\Vert f\Vert_{C^0}(C_{\theta})$.
\end{lemma}

\begin{proof}
    Let
    \[
        u=l^{-1}h.
    \]
    Since $h$ solves Eq. \eqref{Eq}, the function $u$ satisfies
    \begin{equation}\label{a2}
    \left\{
    \begin{aligned}
    &\det\left(l\nabla^2 u+\cos\theta\bigl(\nabla u\otimes e^T
    +e^T\otimes\nabla u\bigr)+u\sigma\right)\\
    &\qquad\qquad
    =f(ul)^{p-1}\left((ul)^2+|\nabla(ul)|^2\right)^{\frac{n+1-q}{2}}
    &&\text{in }C_{\theta},\\
    &\nabla_{\mu}u=0&&\text{on }\partial C_{\theta},
    \end{aligned}
    \right.
    \end{equation}
    where $e^T$ is the tangential part of $e$ on $C_{\theta}$, and we used that $\nabla_{\mu} l = \cot\theta\cdot l$ on $\partial C_{\theta}$.
    Let $\beta = n+1-q$ and  $\zeta=\xi-\cos\theta e$. Then $|\zeta|=1$ on $C_{\theta}$ and
    \[
        l=\sin^2\theta+\cos\theta\langle\xi,e\rangle
        =1+\cos\theta\langle\zeta,e\rangle.
    \]
    Since $-1\leq\langle\zeta,e\rangle\leq-\cos\theta$ on $C_{\theta}$, we have
    \[
        \min\{1-\cos\theta, \sin^2\theta\}\leq l\leq \max\{1-\cos\theta, \sin^2\theta\}.
    \]
    Moreover,
    \[
        |\nabla\langle\zeta,e\rangle|^2=1-\langle\zeta,e\rangle^2,
    \]
    and let $\psi = l^2 + |\nabla l|^2$, therefore
    \[
        \psi
        =(1+\cos\theta\langle\zeta,e\rangle)^2
        +\cos^2\theta\bigl(1-\langle\zeta,e\rangle^2\bigr)
        =2l-\sin^2\theta,
    \]
    hence
    \[
      \min\{  (1-\cos\theta)^2, \sin^2\theta\} \leq \psi\leq \max\{(1-\cos\theta)^2, \sin^2\theta\}.
    \]
    Since $h=ul$ and $\nabla^2l+l\sigma=\sigma$, we have
    \[
        A_{ij}=l u_{ij}+u\sigma_{ij}+l_i u_j+l_j u_i.
    \]
    We also have $\nabla_{\mu}h=\cot\theta\cdot h$ and $\nabla_{\mu}l=\cot\theta\cdot l$ on
    $\partial C_{\theta}$, so
    \[
        \nabla_{\mu}u=0\qquad\text{on }\partial C_{\theta}.
    \]

    At any interior minimum (or maximum) point of $u$ we have $\nabla u=0$ and $\nabla^2u\geq 0$ (or $\nabla^2 u\leq 0$). If an minimum (or maximum) point occurs at the boundary, then
    the tangential derivatives vanish and $\nabla_{\mu}u=0$, we also have $u_{\mu k} = 0$, where $k=1,2,\cdots, n-1$.
    The second derivative in the double normal direction gives
    $u_{\mu\mu}\leq0$ at a boundary maximum point and $u_{\mu\mu}\geq0$ at a boundary
    minimum point. Thus the same Hessian semidefinite conclusion holds at boundary
    extrema point.

    At an extremum point of $u$ the terms $l_i u_j+l_j u_i$ vanish, and
    $
        h^2+|\nabla h|^2=u^2\bigl(l^2+|\nabla l|^2\bigr)=u^2\psi,
    $
    so the equation \eqref{Eq} therefore becomes
    \[
        \det(l\nabla^2u+u\sigma)
        =f u^{n+p-q}l^{p-1}\psi^{\beta/2}.
    \]

    Suppose $u$ attains the maximum value at some point $\xi_0\in C_{\theta}$. At $\xi_0$, we obtain
    \[
        f u^{n+p-q}l^{p-1}\psi^{\beta/2}
        =\det(l\nabla^2u+u\sigma)\leq u^n,
    \]
    which gives
    \[
        u(\xi_0)^{p-q}\leq
        f(\xi_0)^{-1}l(\xi_0)^{1-p}\psi(\xi_0)^{-\beta/2},
    \]
    denote $G:=f^{-1}l^{1-p}\psi^{-\beta/2}$, then 
    $$u(\xi_0)^{p-q} \leq \max_{C_{\theta}}G=:C_1,$$
    where the positive constant $C_1$ depends on $n$, $p$, $q$, $\theta$, $\min_{C_{\theta}}f$ and $\Vert f\Vert_{C^0}(C_{\theta})$.
    Therefore, for every $\xi\in C_{\theta}$, we have
    \[
        h(\xi)^{p-q}
        =u(\xi)^{p-q}l(\xi)^{p-q}
        \leq \max\{1-\cos\theta, \sin^2\theta\}^{p-q}\cdot C_1.
    \]

    Similarly, if $u$ attains the minimum value at $\xi_1\in C_{\theta}$, there holds
    \[
        f u^{n+p-q}l^{p-1}\psi^{\beta/2}
        =\det(l\nabla^2u+u\sigma)\geq u^n,
    \]
    which gives
    \[
        u(\xi_1)^{p-q}\geq
        f(\xi_1)^{-1}l(\xi_1)^{1-p}\psi(\xi_1)^{-\beta/2}
        \geq \min_{C_{\theta}}G=:C_2,
    \]
    where the positive constant $C_2$ depends on $n$, $p$, $q$, $\theta$ and $\Vert f\Vert_{C^0}(C_{\theta})$.
   Therefore, for every $\xi\in C_{\theta}$, we have
    \[
        h(\xi)^{p-q}
        =u(\xi)^{p-q}l(\xi)^{p-q}
        \geq \min\{1-\cos\theta, \sin^2\theta\}^{p-q}\cdot C_2.
    \]
   This completes the proof.
\end{proof}

Next, we give the $C^1$-estimate, which is based solely on convexity and is independent of the specific equation.

\begin{lemma}[{\bf $C^1$ estimate}]\label{c1-1}
    Let $p>q$ and $\theta\in (0,\pi)$. Suppose $h$ is a positive smooth strictly convex solution of Eq. \eqref{Eq}, then there holds
    \begin{equation}\label{b1}
        \max_{C_{\theta}}|\nabla h| \leq (1+\cot^2\theta)^{1/2} \Vert h\Vert_{C^{0}(C_{\theta})}.
    \end{equation}
\end{lemma}

\begin{proof}
    We consider the function
    $$P:= |\nabla h|^2 + h^2.$$
    Suppose that the function $P$  attains its maximum value at some point $\xi_{0}\in C_{\theta}$. If $\xi_0 \in C_{\theta}\backslash\partial C_{\theta}$, we have
    $$0=\nabla_{e_i}P = 2h_{k}h_{ki} + 2hh_{i}, \quad for \quad 1\leq i\leq n.$$
Together with the convexity of $h$, i.e., $h_{ij}+h\delta_{ij}>0$, it follows $\nabla h(\xi_{0})=0$, so \eqref{b1} holds.

If $\xi_{0}\in \partial C_{\theta}$, we choose an orthonormal frame $\{e_i\}_{i=1}^{n}$ around $\xi_{0}\in\partial C_{\theta}$ such that $e_n = \mu$. From Proposition \ref{pro}, we know that
\begin{equation}\label{b2}
    h_{kn}=0 \qquad for~~any \quad 1\leq k\leq n-1.
\end{equation}
Then we have
$$0=\nabla_{e_k}P = 2\sum_{i=1}^{n}h_{i}h_{ik}+ 2hh_{k},$$
which implies
\begin{equation}\label{b3}
    h_{k}(\xi_{0}) = 0.
\end{equation}
Together with \eqref{b3} and $h_n = \cot\theta\cdot h$ on $\partial C_{\theta}$, we have
$$|\nabla h|^2(\xi_{0})\leq (|\nabla h|^2 + h^2)(\xi_0) = (h_n^2 + h^2)(\xi_{0}) = (1+\cot^2\theta)h^2(\xi_0).$$
In conclusion, \eqref{b1} holds.
\end{proof}



\section{$C^2$ estimates} \label{se4}
Now, we establish the priori $C^2$ estimate for the positive solution of Eq. \eqref{Eq}. All covariant
derivatives are taken with respect to the unit spherical metric on $C_{\theta}$.
Roman indices range from $1$ to $n$. At a boundary point, Greek indices range
over the tangential directions $1,\ldots,n-1$, and $e_n=\mu$ denotes the outward
unit normal of $\partial C_{\theta}\subset C_{\theta}$. We keep the notation
$A_{ij}=h_{ij}+h\sigma_{ij}$ and $P=h^2+|\nabla h|^2$ from the preceding
sections, recall $\beta=n+1-q$, and let $ S=\tr(A^{ij})$, where $A^{ij}$ is the inverse of $A_{ij}$.
When $A_{ij}$ is diagonal in an orthonormal frame, its eigenvalues are denoted by
$\lambda_i>0$, we also write $\sigma_k(A)$ for the $k$-th elementary symmetric
function of these eigenvalues. Denote
\[
       F(A):= \log\det A=\tilde f,\qquad
        \tilde f:=(p-1)\log h+\frac{\beta}{2}\log P+\log f.
\]

\begin{theorem}[{\bf{$C^2$ estimate}}]\label{c2 estimate}
Let $n\geq 3$, $p>q$, and $\theta\in (0,\frac{\pi}{2})$. Suppose that $h$ is a
positive $C^4$ strictly convex solution to Eq. \eqref{Eq}. Assume that
\[
        0<m\leq h\leq M,\qquad |\nabla h|\leq K,
\]
then there holds 
\begin{equation}\label{d1}
        \max_{C_{\theta}}|\nabla^2h|\leq C,
\end{equation}
where the positive constant $C$ depending only on
$n,p,q,\theta,m,M,K,\min_{C_{\theta}}f$ and
$\Vert f\Vert_{C^2(C_{\theta})}$, 
\end{theorem}

\begin{proof}
The equation and the lower-order bounds give determinant bounds
\[
        D_-\leq\det A\leq D_+,
\]
where
\[
        D_- =
        \min_{C_{\theta}}f\cdot\min\{m^{p-1},M^{p-1}\}
        \cdot\min\{(m^2)^{\beta/2},(M^2+K^2)^{\beta/2}\}
\]
and
\[
        D_+ =
        \max_{C_{\theta}}f\cdot\max\{m^{p-1},M^{p-1}\}
        \cdot\max\{(m^2)^{\beta/2},(M^2+K^2)^{\beta/2}\}.
\]
Let
\begin{equation}\label{auxiliary Q}
        Q:=\log\sigma_1(A)+ah+b|\nabla h|^2+\log\sigma_{n-1}(A),
\end{equation}
where $b>0$ and $a<0$ will be chosen later. Suppose that $Q$ attains its maximum value at some point $\xi_0\in C_{\theta}$. We divide this proof into two cases: either $\xi_0\in \partial C_{\theta}$ or $\xi_0\in C_{\theta} \backslash \partial C_{\theta}$.

\noindent\textbf{Case 1. $\xi_0\in C_{\theta} \backslash \partial C_{\theta}$.}
In this case, we choose an orthonormal frame $\{e_i\}_{i=1}^n$ around $\xi_0$, such that $A_{ij}=(h_{ij}+h\delta_{ij})$ is diagonal, hence $F^{ij}:=\frac{\partial F(A)}{\partial A_{ij}}$ and $h_{ij}$ are also diagonal at $\xi_0$. In this case, we have
\[
        A_{ij,k}=A_{ik,j},\qquad
        A_{ii,jj}=A_{jj,ii}+A_{ii}-A_{jj}.
\]

Since
\[
        \sigma_{n-1}(A)=\det(A)\,S,
\]
define
\[
        W_i=\frac{1}{\sigma_1(A)}+\lambda_i^{-1}
        -\frac{\lambda_i^{-2}}{S}.
\]
Then at $\xi_0$, we have
\begin{equation}\label{Q stationary}
        0=Q_k
        =\sum_iW_iA_{ii,k}+\bigl(a+2b(\lambda_k-h)\bigr)h_k.
\end{equation}

We need the quadratic third-derivative terms that occur after contracting
second derivatives with $A^{ij}$.

\begin{lemma}\label{third forms}
At a diagonal point where $A_{ij,k}$ is symmetric in all three indices, set
\[
\begin{split}
        \mathcal G_{\sigma_1}
        &=\frac{1}{\sigma_1(A)}
        \sum_{i,j,k}\lambda_i^{-1}\lambda_j^{-1}A_{ij,k}^2
        -\frac{1}{\sigma_1(A)^2}
        \sum_k\lambda_k^{-1}\left(\sum_iA_{ii,k}\right)^2,
\end{split}
\]
and
\[
\begin{split}
        \mathcal G_S
        &=\frac{1}{S}
        \sum_{i,j,k}
        \left(2\lambda_k^{-1}\lambda_i^{-2}\lambda_j^{-1}
        -\lambda_k^{-2}\lambda_i^{-1}\lambda_j^{-1}\right)A_{ij,k}^2\\
        &\qquad
        -\frac{1}{S^2}
        \sum_k\lambda_k^{-1}
        \left(\sum_i\lambda_i^{-2}A_{ii,k}\right)^2.
\end{split}
\]
Then $\mathcal G_{\sigma_1}\geq0$ and $\mathcal G_S\geq0$.
\end{lemma}

\noindent\emph{Proof of Lemma \ref{third forms}.}
Fix the derivative index $k$ and denote $v_i=A_{ii,k}$. For
$\sigma_1(A)\cdot\mathcal G_{\sigma_1}$, the diagonal coefficients after grouping all terms
forced by the three-index symmetry are
\[
        d_k=\lambda_k^{-2},\qquad
        d_i=\lambda_i^{-2}+2\lambda_i^{-1}\lambda_k^{-1}\quad (i\ne k),
\]
and the negative term is
\[
        -\frac{\lambda_k^{-1}}{\sigma_1(A)}
        \left(\sum_i v_i\right)^2.
\]
Cauchy-Schwarz inequality gives
\[
        \left(\sum_i v_i\right)^2
        \leq \left(\sum_i d_i v_i^2\right)\left(\sum_i d_i^{-1}\right).
\]
Here $d_k^{-1}=\lambda_k^2$, and for $i\ne k$,
\[
        d_i^{-1}
        =\frac{\lambda_i^2\lambda_k}{\lambda_k+2\lambda_i}
        \leq \lambda_i\lambda_k.
\]
Thus
\[
        \sum_i d_i^{-1}\leq \lambda_k\sum_i\lambda_i
        =\lambda_k\sigma_1(A),
\]
which makes the grouped diagonal part nonnegative. If $i,j,k$ are pairwise
distinct, the total coefficient of the six symmetric permutations is
\[
        2(\lambda_i^{-1}\lambda_j^{-1}
        +\lambda_i^{-1}\lambda_k^{-1}
        +\lambda_j^{-1}\lambda_k^{-1})\geq0.
\]
Therefore $\mathcal G_{\sigma_1}\geq0$.

Similarly, for $S\cdot\mathcal G_S$, again fix $k$ and write $v_i=A_{ii,k}$. The diagonal
coefficients are
\[
        d_k=\lambda_k^{-4},\qquad
        d_i=\lambda_i^{-2}\lambda_k^{-1}
        (2\lambda_i^{-1}+\lambda_k^{-1})\quad (i\ne k),
\]
and the negative term is
\[
        -\frac{\lambda_k^{-1}}{S}
        \left(\sum_i\lambda_i^{-2}v_i\right)^2.
\]
Cauchy-Schwarz inequality gives
\[
        \left(\sum_i\lambda_i^{-2}v_i\right)^2
        \leq \left(\sum_i d_i v_i^2\right)
        \left(\sum_i\frac{\lambda_i^{-4}}{d_i}\right).
\]
It is enough to prove
\[
        \lambda_k^{-1}
        +\sum_{i\ne k}
        \frac{\lambda_i^{-2}}{2\lambda_i^{-1}+\lambda_k^{-1}}
        \leq S,
\]
and this follows from
\[
        \frac{\lambda_i^{-2}}{2\lambda_i^{-1}+\lambda_k^{-1}}
        \leq \lambda_i^{-1}.
\]
For pairwise distinct $i,j,k$, the coefficient of the six symmetric
permutations is
\[
        2\lambda_i^{-1}\lambda_j^{-1}\lambda_k^{-1}
        (\lambda_i^{-1}+\lambda_j^{-1}+\lambda_k^{-1})\geq0.
\]
Hence $\mathcal G_S\geq0$. This proves the local lemma. \hfill$\Box$

Differentiating $\log\det A=\tilde f$ twice we have
\begin{equation}\label{s4-1}
        \tilde f_{ii}
        =\sum_k\lambda_k^{-1}A_{kk,ii}
        -\sum_{k,s}\lambda_k^{-1}\lambda_s^{-1}A_{ks,i}^2,
\end{equation}
so, for $\log\sigma_1(A)$ we know
\[
        \sum_i\lambda_i^{-1}(\log\sigma_1(A))_{ii}
        =\mathcal G_{\sigma_1}+\frac{1}{\sigma_1(A)}\sum_i\tilde f_{ii}
        +S-\frac{n^2}{\sigma_1(A)}.
\]
For the inverse term $S$, by directly calculate
\[
        S_k=-\sum_i\lambda_i^{-2}A_{ii,k},
\]
and
\begin{equation}\label{s4-2}
        S_{kk}
        =-\sum_i\lambda_i^{-2}A_{ii,kk}
        +2\sum_{i,j}\lambda_i^{-2}\lambda_j^{-1}A_{ij,k}^2.
\end{equation}
Using the formula $A_{ii,jj}=A_{jj,ii}+A_{ii}-A_{jj}$ and \eqref{s4-1}, we have
\[
        \sum_k\lambda_k^{-1}A_{ii,kk}
        =\tilde f_{ii}
        +\sum_{k,s}\lambda_k^{-1}\lambda_s^{-1}A_{ks,i}^2
        +\lambda_i\cdot S-n.
\]
Then the formula \eqref{s4-2} can be written as
\[
\begin{split}
        \sum_k\lambda_k^{-1}S_{kk}
        &=-\sum_i\lambda_i^{-2}\tilde f_{ii}
        -S^2+n\sum_i\lambda_i^{-2}\\
        &\qquad
        +\sum_{i,j,k}
        \left(2\lambda_k^{-1}\lambda_i^{-2}\lambda_j^{-1}
        -\lambda_k^{-2}\lambda_i^{-1}\lambda_j^{-1}\right)A_{ij,k}^2.
\end{split}
\]
After subtracting $S^{-2}\sum_k\lambda_k^{-1}
(\sum_i\lambda_i^{-2}A_{ii,k})^2$, this yields
\begin{equation}\label{contracted log S}
        \sum_k\lambda_k^{-1}(\log S)_{kk}
        =\mathcal G_S-\frac{1}{S}\sum_i\lambda_i^{-2}\tilde f_{ii}
        -S+\frac{n\sum_i\lambda_i^{-2}}{S}.
\end{equation}
Also
\[
        \sum_i\lambda_i^{-1}h_{ii}=n-hS
\]
and
\[
\begin{split}
        \sum_i\lambda_i^{-1}(|\nabla h|^2)_{ii}
        &=2\sum_kh_k\tilde f_k+2\sigma_1(A)-4nh
        +2h^2S-2\sum_i\lambda_i^{-1}h_i^2.
\end{split}
\]
Combining the preceding formulas with $\log\sigma_{n-1}(A)=\tilde f+\log S$, at $\xi_0$,
we get
\begin{equation}\label{contracted Q}
\begin{split}
        0&\geq\sum_i\lambda_i^{-1}Q_{ii}\\
        &=\mathcal G_{\sigma_1}+\mathcal G_S+\sum_iW_i\tilde f_{ii}
        -\frac{n^2}{\sigma_1(A)}+a(n-hS)\\
        &\qquad
        +b\left(2\sum_kh_k\tilde f_k+2\sigma_1(A)-4nh
        +2h^2S-2\sum_i\lambda_i^{-1}h_i^2\right)
        +\frac{n\sum_i\lambda_i^{-2}}{S}.
\end{split}
\end{equation}
Lemma \ref{third forms} makes the first two terms nonnegative, and easy find that the last term is also
nonnegative.

Since $\tilde f:=(p-1)\log h+\frac{\beta}{2}\log P+\log f$, by directly calculate we know
\begin{equation}\label{f second derivatives}
\begin{split}
        \tilde f_{ii}
        &=\beta\left(
        \frac{\lambda_i^2}{P}
        -\frac{2h_i^2\lambda_i^2}{P^2}
        +\frac{\sum_k h_kA_{ii,k}}{P}
        -\frac{h\lambda_i}{P}\right)\\
        &\qquad
        +(p-1)\left(\frac{\lambda_i}{h}-1-\frac{h_i^2}{h^2}\right)
        +(\log f)_{ii}.
\end{split}
\end{equation}
Using $Q_k=0$, then
\[
      \frac{\beta}{P}
        \sum_kh_k\sum_iW_iA_{ii,k}=  -\frac{a\beta|\nabla h|^2}{P}
        -\frac{2b\beta}{P}\sum_kh_k^2\lambda_k
        +\frac{2b\beta h|\nabla h|^2}{P}.
\]
On the other hand,
\[
        2b\sum_kh_k\tilde f_k
        =2b(p-1)\frac{|\nabla h|^2}{h}
        +\frac{2b\beta}{P}\sum_kh_k^2\lambda_k
        +2b\sum_kh_k(\log f)_k.
\]
The terms containing $\sum_kh_k^2\lambda_k$ cancel.

Since
$P\geq m^2$ and $|h_i|^2\leq K^2$, the $\lambda_i^2$ terms above
are bounded from below by $-C\sum_iW_i\lambda_i^2$. For $W_i$ we have
\[
        0\leq W_i,\qquad
        \sum_iW_i\lambda_i^2\leq 2\sigma_1(A),\qquad
        \sum_iW_i\lambda_i= n,\qquad
        \sum_iW_i\leq C+S,
\]
where we used $\sigma_1(A)\geq nD_-^{1/n}$. Thus the terms containing
$\beta$ and $p-1$ are controlled for all signs of these numbers. The
$(\log f)_{ii}$ terms are bounded below by some negative constant $-C$, and
\[
        2h^2S-2\sum_i\lambda_i^{-1}h_i^2\geq -2K^2S.
\]

So the inequality \eqref{contracted Q} can be written as
\begin{equation}\label{interior inequality}
\begin{split}
        0\geq
        (2b-C)\sigma_1(A)+(-am-2bK^2-C)S
        -C(1+|a|+b).
\end{split}
\end{equation}
Choose $b=b_1>C+1$. Then choose $a<a_1<0$ such that
\[
        -a_1m-2bK^2-C\geq 1.
\]
Then at $\xi_0$, the formula \eqref{interior inequality} can be written as 
\[
        0\geq \sigma_1(A) - C(1+|a|+b).
\]
So we conclude that $\sigma_1(\xi_0)\leq C$, where the positive constant $C$ depends on $n, p, q, \min_{C_{\theta}}f, \min_{C_{\theta}}h, \Vert f\Vert_{C^2 (C_{\theta})}$ and $\Vert h\Vert _{C^0 (C_{\theta})}$.

\noindent\textbf{Case 2. $\xi_0\in \partial C_{\theta}$}
We choose an orthonormal frame $\{e_i\}_{i=1}^n$ around $\xi_0\in\partial C_{\theta}$ satisfying $e_n=\mu$ at $\xi_0$. So we know 
\[
        \lambda_n=A_{nn},\qquad
        \lambda_{\alpha}=A_{\alpha\alpha}, \quad \alpha\in 1,2,\cdots, n-1
\]

At the boundary, with Robin condition $h_n=\cot\theta\,h$, we know that $h_{\alpha n}=0$.
By directly calculate, we know the normal derivative of $\tilde{f}$ is 
\[
        \tilde f_n
        =(p-1)\cot\theta
        +\beta\cot\theta\,
        \frac{h\lambda_n}{P}
        +(\log f)_n,
\]
by $P_n=2\cot\theta\,h\lambda_n$. Differentiating
$\log\det A=\tilde f$ in the normal direction and using
$S=\lambda_n^{-1}+\sum_{\alpha<n}\lambda_{\alpha}^{-1}$ we have
\[
        (\lambda_n)_n
        =\lambda_n\tilde f_n
        -\cot\theta\,\lambda_n(\lambda_nS-n).
\]
Then
\[
\begin{split}
        (\sigma_1(A))_n
        &=(\lambda_n)_n+\sum_{\alpha<n}(\lambda_{\alpha})_n\\
        &=\lambda_n\tilde f_n
        +\cot\theta\bigl(2n\lambda_n-\sigma_1(A)-\lambda_n^2S\bigr),
\end{split}
\]
and
\[
\begin{split}
        S_n
        &=-\lambda_n^{-2}(\lambda_n)_n
        -\sum_{\alpha<n}\lambda_{\alpha}^{-2}(\lambda_{\alpha})_n\\
        &=-\frac{\tilde f_n}{\lambda_n}
        -\cot\theta\,\lambda_n
        \sum_{\alpha<n}\left(\frac1{\lambda_{\alpha}}
        -\frac1{\lambda_n}\right)^2.
\end{split}
\]
It follows from $\sigma_{n-1}(A)=\det(A)S$ that
\[
        (\log\sigma_{n-1}(A))_n
        =\left(1-\frac1{\lambda_nS}\right)\tilde f_n
        -\frac{\cot\theta\,\lambda_n}{S}
        \sum_{\alpha<n}\left(\frac1{\lambda_{\alpha}}
        -\frac1{\lambda_n}\right)^2.
\]
Together with $|\nabla h|^2_n=2\cot\theta\,h(\lambda_n-h)$ and
\[
        2n-\frac{\sigma_1(A)}{\lambda_n}-\lambda_nS
        =-\sum_{\alpha<n}\left(\frac{\lambda_n}{\lambda_{\alpha}}
        +\frac{\lambda_{\alpha}}{\lambda_n}-2\right). 
\]
Then, at $\xi_0$, we have
\begin{equation}\label{boundary Q derivative}
\begin{split}
        0\leq \frac{\sigma_1(A)}{\lambda_n}Q_n
        &=\left(1+\frac{\sigma_1(A)}{\lambda_n}
        -\frac{\sigma_1(A)}{\lambda_n^2S}\right)\tilde f_n\\
        &\qquad
        -\cot\theta
        \sum_{\alpha<n}\left(\frac{\lambda_n}{\lambda_{\alpha}}
        +\frac{\lambda_{\alpha}}{\lambda_n}-2\right)\\
        &\qquad
        +a\cot\theta\,h\frac{\sigma_1(A)}{\lambda_n}
        +2b\cot\theta\,h(\lambda_n-h)\frac{\sigma_1(A)}{\lambda_n}\\
        &\qquad
        -\cot\theta\,\frac{\sigma_1(A)}{S}
        \sum_{\alpha<n}\left(\frac1{\lambda_{\alpha}}
        -\frac1{\lambda_n}\right)^2.
\end{split}
\end{equation}

First suppose $\lambda_n\leq m/2$. Then $\lambda_n-h\leq -m/2$, so the
 term
 $$2b\cot\theta h(\lambda_n -h)\frac{\sigma_1(A)}{\lambda_n}$$
 is nonpositive. Then we can written \eqref{boundary Q derivative} as

 \begin{equation}\label{boundary Q derivative 2}
\begin{split}
        0&\leq \left(1+\frac{\sigma_1(A)}{\lambda_n}
        -\frac{\sigma_1(A)}{\lambda_n^2S}\right)\tilde f_n
        -\cot\theta
        \sum_{\alpha<n}\left(\frac{\lambda_n}{\lambda_{\alpha}}
        +\frac{\lambda_{\alpha}}{\lambda_n}-2\right)\\
        &\qquad
       +a\cot\theta\,h\frac{\sigma_1(A)}{\lambda_n} -\cot\theta\,\frac{\sigma_1(A)}{S}
        \sum_{\alpha<n}\left(\frac1{\lambda_{\alpha}}
        -\frac1{\lambda_n}\right)^2  \\ 
        & 
        \leq \left(1+\frac{\sigma_1(A)}{\lambda_n}
        -\frac{\sigma_1(A)}{\lambda_n^2S}\right)\tilde f_n
        -\cot\theta
        \sum_{\alpha<n}\left(\frac{\lambda_n}{\lambda_{\alpha}}
        +\frac{\lambda_{\alpha}}{\lambda_n}-2\right) +a\cot\theta\,h\frac{\sigma_1(A)}{\lambda_n}\\ 
        &\leq 
        \left(1+\frac{\sigma_1(A)}{\lambda_n}
        -\frac{\sigma_1(A)}{\lambda_n^2S}\right)\tilde f_n +a\cot\theta\,h\frac{\sigma_1(A)}{\lambda_n},
\end{split}
\end{equation}
 The last inequality holds by using $z+z^{-1}-2 = \frac{(z-1)^2}{z} ~~(z>0)$.
We know that 
$$ \tilde f_n
        =(p-1)\cot\theta
        +\beta\cot\theta\,
        \frac{h\lambda_n}{P}
        +(\log f)_n \leq C+C\lambda_n,
        $$
where the positive constant $C$ depend on $n, p, q, \min_{C_{\theta}}f, \min_{C_{\theta}}h, \Vert f\Vert_{C^2 (C_{\theta})}$ and $\Vert h\Vert _{C^0 (C_{\theta})}$. Since we suppose $\lambda_n \leq m/2$, so we have $\tilde{f}_n\leq C$. So we can rewritten the formula \eqref{boundary Q derivative 2} as 
\begin{equation}\label{s4-3}
0\leq C+\left(C+a\cot\theta\,m\right)
        \frac{\sigma_1(A)}{\lambda_n}.
\end{equation}

Then choose $a<a_2<0$ such that $C+a_2 \cot\theta m<-1$. Then we have
$$0\leq C-\frac{\sigma_1(A)}{\lambda_n}\leq C-\frac{2}{m}\sigma_1(A), $$
i.e.,
$$\sigma_1(A)\leq \frac{mC}{2}.$$

Next we consider the case that $\lambda_n\geq m/2$. In this case we know 
$\tilde f_n\leq C(1+\lambda_n)$, and
\[
        \left(1+\frac{\sigma_1(A)}{\lambda_n}
        -\frac{\sigma_1(A)}{\lambda_n^2S}\right)\tilde f_n
        \leq C(1+\sigma_1(A)).
\]

Then the formula \eqref{boundary Q derivative} can be written as
\begin{equation}\label{boundary inequality} 
\begin{split}
        0&\leq 
        C(1+\sigma_1(A))
        +a\cot\theta\,h\frac{\sigma_1(A)}{\lambda_n}-\cot\theta\,\frac{\sigma_1(A)}{S}
        \sum_{\alpha<n}\left(\frac1{\lambda_{\alpha}}
        -\frac1{\lambda_n}\right)^2\\
        &\leq 
        C(1+\sigma_1(A))-\cot\theta\,\frac{\sigma_1(A)}{S}
        \sum_{\alpha<n}\left(\frac1{\lambda_{\alpha}}
        -\frac1{\lambda_n}\right)^2.
\end{split}
\end{equation}

By Cauchy-Schwarz inequality,
\[
        \sum_{\alpha<n}\left(\frac1{\lambda_{\alpha}}
        -\frac1{\lambda_n}\right)^2
        \geq\frac1{n-1}\left(S-\frac n{\lambda_n}\right)^2.
\]
If $S\geq4n/m$, then $\lambda_n\geq m/2$ implies
$S-n/\lambda_n\geq S/2$. Therefore we have
\begin{equation}\label{s4-4}
        0\leq C(1+\sigma_1(A))
        -\frac{\cot\theta}{4(n-1)}\sigma_1(A)S.
\end{equation}
Since $\sigma_1(A)\geq nD_-^{1/n}>0$, divide both sides by $\sigma_1(A)$, then we have
$$0\leq C - \frac{\cot\theta}{4(n-1)}S,$$
so $S$ has a positive upper bound. If $S<4n/m$, it is already giving such a bound.
Hence $S\leq C$. Each eigenvalue satisfies
$\lambda_i\geq S^{-1}$, so we have
\[
        \lambda_i
        =\frac{\det A}{\prod_{j\ne i}\lambda_j}
        \leq D_+S^{n-1}\leq C.
\]
Thus $\sigma_1(A)(\xi_0)\leq C$. In conclusion, by setting $a=\min\{a_1, a_2\}, b=b_1$, we can obtain $\sigma_1(A)(\xi_0)\leq C$ at the maximum point $\xi_0\in C_{\theta}$ of $Q$. Then at $\xi_0$, we have 
$$\sigma_{n-1}(A)(\xi_0)\leq n\sigma_1(A)(\xi_0)^{n-1}\leq nC^{n-1},$$
then 
$$Q(\xi_0)\leq C.$$
On the other hand, by MacLaurin inequality, we have
$$\sigma_{n-1}(A)\geq n(\det A)^{(n-1)/n}\geq nD^{(n-1)/n}_-=: c_0>0.$$
In conclusion, at any point $\xi\in C_{\theta}$, we have 
$$\log \sigma_1(A)(\xi) + \log c_0 +aM\leq Q(\xi)\leq Q(\xi_0)\leq C,$$
that means $\sigma_1(A)\leq C$ in $C_{\theta}$.
\end{proof}

The curvature estimates imply the regularity bound needed for the continuity method.

\begin{corollary}[]\label{c2 theorem}
Let $n\geq 3$, $p>q$, and $\theta\in (0,\frac{\pi}{2})$. Suppose that $h$ is a
smooth positive strictly convex solution to Eq. \eqref{Eq}. Then
\[
        \min_{C_{\theta}}h\geq c,
\]
and, for any $\alpha\in(0,1)$,
\[
        \Vert h\Vert_{C^{3,\alpha}(C_{\theta})}\leq C,
\]
where the positive constants $c$ and $C$ depend only on
$n,p,q,\theta,\alpha$, and $f$.
\end{corollary}

\begin{proof}
Combining Lemma \ref{C0 estimate}, Lemma \ref{c1-1} and Theorem \ref{c2 estimate}, we obtain
$$c\leq \min_{C_{\theta}}h, \qquad and \qquad \Vert h\Vert_{C^2(C_{\theta})}\leq C.$$
By the theory of fully nonlinear second-order uniformly elliptic equations and the Schauder estimate we can proof this corollary.

\end{proof}

 
\section{Proof of Theorem \ref{main1.1}} \label{se5}
We use the continuity method in logarithmic variables. For a positive function
$h$, set
\[
        v=\log h,\qquad
        B[v]_{ij}=v_{ij}+v_iv_j+\sigma_{ij}.
\]
Then
\[
        h_{ij}+h\sigma_{ij}=e^vB[v]_{ij},\qquad
        h^2+|\nabla h|^2=e^{2v}\bigl(1+|\nabla v|^2\bigr),
\]
and the Robin condition becomes
\[
        \nabla_{\mu}h=\cot\theta\,h
        \quad\Longleftrightarrow\quad
        \nabla_{\mu}v=\cot\theta.
\]

Define
\[
        f^{(0)}
        =l^{1-p}\left(l^2+|\nabla l|^2\right)^{\frac{q-n-1}{2}},
        \qquad
        f_t=(1-t)f^{(0)}+tf,\qquad 0\leq t\leq1.
\]
Consider the following family of equations
\[
\left\{
\begin{aligned}
&\det(\nabla^2h+h\sigma)
=h^{p-1}\left(h^2+|\nabla h|^2\right)^{\frac{n+1-q}{2}}f_t
&&\text{in }C_{\theta},\\
&\nabla_{\mu}h=\cot\theta\,h
&&\text{on }\partial C_{\theta}.
\end{aligned}
\right.
\]
Equivalently,
\begin{equation}\label{log path}
\left\{
\begin{aligned}
&\Phi_t[v]:=
\log\det B[v]
-\frac{n+1-q}{2}\log\bigl(1+|\nabla v|^2\bigr)
-(p-q)v-\log f_t=0
&&\text{in }C_{\theta},\\
&\nabla_{\mu}v=\cot\theta
&&\text{on }\partial C_{\theta}.
\end{aligned}
\right.
\end{equation}

\begin{proof}[Proof of Theorem~\ref{main1.1}]
Let $0<\alpha<1$ and put
\[
        \mathcal V_{\alpha}
        =\{v\in C^{2,\alpha}(C_{\theta}): B[v]>0\}.
\]
Define
\[
        \Psi(t,v)=\bigl(\Phi_t[v],v_{\mu}-\cot\theta\bigr)
\]
as a map from $[0,1]\times\mathcal V_{\alpha}$ to
$C^{0,\alpha}(C_{\theta})\times C^{1,\alpha}(\partial C_{\theta})$, and set
\[
        \mathcal I
        =\{t\in[0,1]:\text{ there is }v\in\mathcal V_{\alpha}
        \text{ with }\Psi(t,v)=0\}.
\]

The set $\mathcal I$ is nonempty. The identities for $l$ used in the proof of
Lemma \ref{C0 estimate},
\[
        \nabla^2l+l\sigma=\sigma,\qquad
        l_{\mu}=\cot\theta\,l,
\]
show that $h=l$ is strictly convex and satisfies the boundary condition. With
the above definition of $f^{(0)}$, it also satisfies the path equation at
$t=0$. Hence $0\in\mathcal I$.

We next prove openness. At an admissible zero $v$, the linearization in a
variation $\phi$ is
\begin{equation}\label{linearized log}
        D_v\Phi_t[v]\phi
        =B^{ij}\phi_{ij}
        +\left(2B^{ij}v_i
        -\frac{(n+1-q)v^j}{1+|\nabla v|^2}\right)\phi_j
        -(p-q)\phi,
\end{equation}
where $(B^{ij})$ is the inverse of $(B[v]_{ij})$; the same convention is used
for $B_a$ below. The boundary linearization is $\phi_{\mu}=0$.
Since $B[v]>0$ on the compact domain, this operator is uniformly elliptic, and the boundary operator is uniformly
oblique. The zeroth-order coefficient is $-(p-q)<0$. The maximum principle and
Hopf boundary lemma \cite[Lemma~3.4 and Theorem~3.5]{GT83} give zero kernel for
the homogeneous oblique problem. The oblique Schauder estimate
\cite[Theorem~6.30]{GT83}, together with the linear solvability theorem
\cite[Theorem~6.31]{GT83}, gives an isomorphism
\[
        \phi\mapsto (D_v\Phi_t[v]\phi,\phi_{\mu})
\]
from $C^{2,\alpha}(C_{\theta})$ to
$C^{0,\alpha}(C_{\theta})\times C^{1,\alpha}(\partial C_{\theta})$. The Banach
implicit-function theorem therefore shows that $\mathcal I$ is open.

Every admissible $C^{2,\alpha}$ solution is smooth up to the
boundary. Indeed, for a fixed such solution, the equation is
uniformly elliptic with $C^{0,\alpha}$ linearized coefficients,
and the boundary condition is linear and uniformly oblique.
The standard bootstrap using tangential difference quotients,
oblique Schauder estimates \cite[Section~6.7]{GT83}, and
recovery of normal derivatives from the equation yields
$v\in C^\infty(C_{\theta})$.

We prove closedness. Since $\theta\in(0,\frac{\pi}{2})$, the function $l$ is
bounded below by a positive constant. Thus $f^{(0)}$ is positive and smooth, and
for every $k$,
\[
        \inf_{C_{\theta}}f_t
        \geq \min\{\inf_{C_{\theta}}f^{(0)},\inf_{C_{\theta}}f\}>0,
        \qquad
        \Vert f_t\Vert_{C^k(C_{\theta})}
        \leq \Vert f^{(0)}\Vert_{C^k(C_{\theta})}
        +\Vert f\Vert_{C^k(C_{\theta})},
\]
uniformly in $t$. Let $t_j\in\mathcal I$ and $t_j\to t_{\infty}$. Choose
solutions $v_j$ and set $h_j=e^{v_j}$. By the preceding bootstrap, each $h_j$ is
smooth. Lemmas \ref{C0 estimate} and \ref{c1-1}, Theorem \ref{c2 estimate}, and
Corollary \ref{c2 theorem}, applied with $f_t$ in place of $f$, give uniform bounds
\[
        0<m\leq h_j\leq M,\qquad
        |\nabla h_j|\leq K,\qquad
        C^{-1}\sigma\leq A[h_j]\leq C\sigma,
        \qquad
        \Vert h_j\Vert_{C^{3,\gamma}(C_{\theta})}\leq C
\]
for some $\gamma\in(\alpha,1)$. After passing to a subsequence,
\[
        h_j\to h_{\infty}\qquad\text{in }C^{2,\alpha}(C_{\theta}).
\]
The limit is positive, the two-sided bound for $A[h_j]$ passes to the limit, and
therefore $A[h_{\infty}]>0$. Passing to the limit in the path equation and in
the boundary condition gives a solution at $t_{\infty}$. Thus
$v_{\infty}=\log h_{\infty}$ lies in $\mathcal V_{\alpha}$ and satisfies
$\Psi(t_{\infty},v_{\infty})=0$. Hence $t_{\infty}\in\mathcal I$, and
$\mathcal I$ is closed.

Since $[0,1]$ is connected, $\mathcal I=[0,1]$. At $t=1$ we obtain a smooth
positive strictly convex solution of Eq. \eqref{Eq}.

\enlargethispage{\baselineskip}
It remains to prove uniqueness. Let $v_0,v_1\in\mathcal V_{\alpha}$ be two
solutions for the same $t$, and set $w=v_1-v_0$. Then $w_{\mu}=0$ on
$\partial C_{\theta}$. Put $B_a=B[v_a]$ for $a=0,1$. Concavity of $\log\det$ on
the positive cone gives
\[
        \log\det B_1-\log\det B_0
        \leq B_0^{ij}(B_1-B_0)_{ij}.
\]
Since
\[
        (B_1-B_0)_{ij}=w_{ij}+v_{1,i}w_j+w_i v_{0,j},
\]
the first-order terms combine after contraction with the symmetric tensor
$B_0^{ij}$. The gradient logarithm has the exact line-segment identity
\[
\begin{split}
        &\log\bigl(1+|\nabla v_1|^2\bigr)
        -\log\bigl(1+|\nabla v_0|^2\bigr)\\
        &\qquad =
        2\int_0^1
        \frac{v_0^j+s w^j}{1+|\nabla v_0+s\nabla w|^2}\,ds\, w_j.
\end{split}
\]
Subtracting the two equations gives the linear inequality
\begin{equation}\label{comparison inequality}
\begin{split}
        0&\leq B_0^{ij}w_{ij}-(p-q)w\\
        &\quad+\left(B_0^{ij}(v_{1,i}+v_{0,i})
        -(n+1-q)\int_0^1
        \frac{v_0^j+s w^j}{1+|\nabla v_0+s\nabla w|^2}\,ds\right)w_j.
\end{split}
\end{equation}
The operator is uniformly elliptic because $B_0>0$ on the compact cap. A
positive interior maximum of $w$ contradicts the displayed inequality. A positive
boundary maximum contradicts the same inequality together with the Hopf boundary
lemma and the condition $w_{\mu}=0$. Thus $w\leq0$. Reversing the roles of
$v_0$ and $v_1$ gives $w\geq0$, hence $v_0\equiv v_1$ and uniqueness follows.
\end{proof}

%\section*{Acknowledgments}
% Acknowledgments and funding information will be supplied later.

\clearpage
\begin{thebibliography}{50}
\setlength{\itemsep}{-0pt} \small

\bibitem{Ale38} A. D. Aleksandrov, \emph{On the theory of mixed volumes. III. Extensions of two theorems of Minkowski on convex polyhedra to arbitrary convex bodies}, Mat. Sb., {\bf 3} (1938) 27-46.

\bibitem{Ale39} A. D. Aleksandrov, \emph{On the surface area measure of convex bodies}, Mat. Sb., {\bf 6} (1939) 167-174.

\bibitem{Ale56} A. D. Aleksandrov, \emph{Uniqueness theorems for surfaces in the large. I}, Vestnik Leningrad. Univ., 11.19 (1956) 5-17.

\bibitem{BF.JDE.266-2019.7980} K. J. B\"{o}r\"{o}czky and F. Fodor, \emph{The $L_p$ dual Minkowski problem for $p>1$ and $q>0$}, J. Differential Equations, {\bf266} (2019) 7980-8033.

\bibitem{BHP.JDG.109-2018.411}
K. J. B\"{o}r\"{o}czky, M. Henk and H. Pollehn, \emph{Subspace
  concentration of dual curvature measures of symmetric convex bodies}, J. Differential Geom., {\bf 109} (2018) 411-429.

\bibitem{BLYZ.JAMS.26-2013.831}
K. J. B{\"o}r{\"o}czky, E. Lutwak, D. Yang and G. Zhang, \emph{The
  logarithmic Minkowski problem}, J. Amer. Math. Soc., {\bf 26} (2013) 831-852.

  \bibitem{Busemann59} H. Busemann, \emph{Minkowski's and related problems for convex surfaces with boundaries}, Mich. Math. J., {\bf 6} (1959) 259-266.
  
  \bibitem{CHZ.MA.373-2019.953}
C. Chen, Y. Huang and Y. Zhao, \emph{Smooth solutions to the $L_p$ dual Minkowski problem}, Math. Ann., {\bf 373} (2019) 953-976.

\bibitem{CCL} H. Chen, S. Chen and Q. R. Li, \emph{Variations of a class of Monge-Amp\`ere type functionals and their applications}, Anal. PDE., {\bf 14} (2021) 689-716.

\bibitem{CLZ.TAMS.371-2019.2623} S. Chen, Q. R. Li and G. Zhu, \emph{The logarithmic Minkowski problem for non-symmetric measures}, Trans. Amer. Math. Soc., {\bf 371} (2019) 2623-2641.

\bibitem{CY76} S. Y. Cheng and S. T. Yau, \emph{On the regularity of the solution of the $n$-dimensional Minkowski problem}, Comm. Pure Appl. Math., {\bf 29} (1976) 495-516.

\bibitem{CW95} K. S. Chou, X. J. Wang, \emph{Minkowski problems for complete noncompact convex hypersurfaces}, Topol. Methods Nonlinear Anal., {\bf 6} (1995) 151-162.

\bibitem{CW.Adv.205-2006.33}
K. S. Chou and X. J. Wang, \emph{The $L_p$-Minkowski problem and the Minkowski problem in centroaffine geometry}, Adv. Math., {\bf 205} (2006) 33-83.

\bibitem{GY26v2} Y. Gao, \emph{The $L_p$ dual Minkowski problem for capillary hypersurfaces}, arXiv: 2510.12804v2.

\bibitem{GT83} D. Gilbarg, N. Trudinger, \emph{Elliptic Partial Differential Equations of Second Order}, second edition, Grundlehren der Mathematischen Wissenschaften (Fundamental Principles of Mathematical Sciences), vol. 224, Springer-Verlag, Berlin, 1983. 


\bibitem{HP.Adv.323-2018.114}
M. Henk and H. Pollehn, \emph{Necessary subspace concentration conditions for the even dual Minkowski problem}, Adv. Math., {\bf 323} (2018) 114-141.

\bibitem{HI25} Y. X. Hu and M. N. Ivaki, \emph{Capillary curvature images}, arXiv: 2505.12921.

\bibitem{HHI25} J. R. Hu, Y. X. Hu and M. N. Ivaki, \emph{Capillary $L_p$ Minkowski flows}, arXiv: 2509.06110.

\bibitem{HIS25} Y. X. Hu, M. N. Ivaki and J. Scheuer, \emph{Capillary Christoffel-Minkowski problem}, arXiv:2504.09320.

\bibitem{HY26} J. R. Hu and B. Yang, \emph{Capillary John ellipsoid theorem with applications to capillary curvature problems}, arXiv:2603.27252v2, 2026.

\bibitem{HJ.JFA.277-2019.2209}
Y. Huang and Y. Jiang, \emph{Variational characterization for the planar dual Minkowski problem}, J. Funct. Anal., {\bf 277} (2019) 2209-2236.

\bibitem{HLYZ.Acta.216-2016.325}
Y. Huang, E. Lutwak, D. Yang and G. Zhang, \emph{Geometric measures in the dual Brunn-Minkowski theory and their associated Minkowski
  problems}, Acta Math., {\bf 216} (2016) 325-388.

\bibitem{HLYZ.JDG.110-2018.1}
\leavevmode\vrule height 2pt depth -1.6pt width 23pt, \emph{The $L_p$-Aleksandrov problem for $L_p$-integral curvature}, J. Differential Geom., {\bf 110} (2018) 1-29.

\bibitem{HZ.Adv.332-2018.57}
Y. Huang and Y. Zhao, \emph{On the $L_p$ dual Minkowski problem}, Adv. Math., {\bf 332} (2018) 57-84.

\bibitem{HLYZ.DCG.33-2005.699}
D. Hug, E. Lutwak, D. Yang and G. Zhang, \emph{On the $L_p$ Minkowski problem for polytopes}, Discrete Comput. Geom., {\bf 33} (2005) 699-715.

\bibitem{JLW.JFA.274-2018.826}
H. Jian, J. Lu and X. J. Wang, \emph{A priori estimates and existence of solutions to the prescribed centroaffine curvature problem}, J. Funct. Anal., {\bf 274} (2018) 826-862.

\bibitem{JLZ.CVPDE.55-2016.41}
H. Jian, J. Lu and G. Zhu, \emph{Mirror symmetric solutions to the
  centro-affine Minkowski problem}, Calc. Var. Partial Differential
  Equations, {\bf 55} (2016) Art. 41, 22 pp.

\bibitem{Lew38} H. Lewy, \emph{On differential geometry in the large. I. Minkowski's problem}, Trans. Amer. Math. Soc., {\bf 43} (1938) 258-270.

\bibitem{LLL}
Q. R. Li, J. Liu and J. Lu, \emph{Non-uniqueness of solutions to the dual $L_p$-Minkowski problem}, IMRN, {\bf 12} (2022) 9114-9150.

\bibitem{LSW.JEMSJ.22-2020.893}
Q. R. Li, W. Sheng and X. J. Wang, \emph{Flow by Gauss curvature to the Aleksandrov and dual Minkowski problems}, J. Eur. Math. Soc. (JEMS), {\bf 22} (2020) 893-923.

%\bibitem{LT86} G. M. Lieberman and N. S. Trudinger, \emph{Nonlinear oblique boundary value problems for nonlinear elliptic equations}, Trans. Amer. Math. Soc., {\bf 295} (1986) 509--546.

\bibitem{LTU86}
P. L. Lions, N. S. Trudinger and J. I. E. Urbas, \emph{The Neumann problem for equations of Monge-Amp\`ere type}, Comm. Pure Appl. Math., {\bf 39} (1986) 539-563.

\bibitem{Lu.SCM.61-2018.511}
J. Lu, \emph{Nonexistence of maximizers for the functional of the
  centroaffine Minkowski problem}, Sci. China Math., {\bf 61} (2018) 511-516.

\bibitem{Lu.JDE.266-2019.4394}
\leavevmode\vrule height 2pt depth -1.6pt width 23pt, \emph{A remark on rotationally symmetric solutions to the centroaffine Minkowski problem}, J. Differential Equations, {\bf 266} (2019) 4394-4431.

\bibitem{LW.JDE.254-2013.983}
J. Lu and X. J. Wang, \emph{Rotationally symmetric solutions to the $L_p$-Minkowski problem}, J. Differential Equations, {\bf 254} (2013) 983-1005.

\bibitem{Lut.JDG.38-1993.131}
E. Lutwak, \emph{The Brunn-Minkowski-Firey theory. I. Mixed volumes and the Minkowski problem}, J. Differential Geom., {\bf 38} (1993) 131-150.

\bibitem{Lut.Adv.118-1996.244}
E. Lutwak, \emph{The Brunn-Minkowski-Firey theory. II. Affine and geominimal surface areas}, Adv. Math., {\bf 118} (1996) 244-294.

\bibitem{LYZ.TAMS.356-2004.4359}
E. Lutwak, D. Yang and G. Zhang, \emph{On the $L_p$-Minkowski
  problem}, Trans. Amer. Math. Soc., {\bf 356} (2004) 4359-4370.

\bibitem{LYZ18} E. Lutwak, D. Yang and G. Zhang, \emph{$L_p$ dual curvature measures}, Adv. Math., {\bf 329} (2018) 85-132.


\bibitem{MWW25a} X. Q. Mei, G. F. Wang and L. J. Weng, \emph{The capillary Minkowski problem}, Adv. Math., {\bf 469} (2025) 110230.

\bibitem{MWW25b} 
\leavevmode\vrule height 2pt depth -1.6pt width 23pt, \emph{The capillary $L_p$-Minkowski problem}, arXiv: 2505.07746.

\bibitem{MWW25c}
\leavevmode\vrule height 2pt depth -1.6pt width 23pt, \emph{The capillary Gauss curvature flow}, arXiv:2506.09840.


\bibitem{Min97} H. Minkowski, \emph{Allgemeine Lehrs\"atze \"uber die konvexen Polyeder}, Nachr. Ges. Wiss. Gottingen, (1897) 198-219.

\bibitem{Min03} H. Minkowski, \emph{Volumen und Oberfl\"ache}, Math. Ann., {\bf 57} (1903) 447-495.

\bibitem{Nir53} L. Nirenberg, \emph{The Weyl and Minkowski problems in differential geometry in the large}, Comm. Pure Appl. Math., {\bf 6} (1953) 337-394.

\bibitem{Oli82} V. Oliker, \emph{The boundary value Minkowski problem. The parametric case}, Ann. Sc. Norm. Super. Pisa, Cl. Sci., {\bf 9} (1982) 463-490.

\bibitem{Pog52} A. V. Pogorelov, \emph{Regularity of a convex surface with given Gaussian curvature}, Mat. Sbornik N.S., {\bf 31(73)} (1952) 88-103.

\bibitem{Pog78} A. V. Pogorelov, \emph{The Minkowski multidimensional problem}, in: Scripta Series in Mathematics, Halsted Press [John Wiley and Sons], New York-Toronto-London, 1978, p.106, Translated from the Russian by Vladimir Oliker, Introduction by Louis Nirenberg. V. H. Winston and Sons, Washington, DC.

\bibitem{Sch18} R. Schneider, \emph{A Brunn-Minkowski theory for coconvex sets of finite volume}, Adv. Math., {\bf 332} (2018) 199-234.

\bibitem{Sch03} O. C. Schn\"urer, \emph{A generalized Minkowski problem with Dirichlet boundary condition}, Trans. Amer. Math. Soc., {\bf 355} (2003) 655-663.

\bibitem{Sta.Adv.167-2002.160}
A. Stancu, \emph{The discrete planar $L_0$-Minkowski problem}, Adv.
  Math., {\bf 167} (2002) 160-174.

\bibitem{WZ25}
X. D. Wang and B. C. Zhu, \emph{The capillary Orlicz-Minkowski problem}, arXiv:2509.10859.

\bibitem{Zha.CVPDE.56-2017.18}
Y.~Zhao, \emph{The dual Minkowski problem for negative indices}, Calc. Var. Partial Differential Equations, {\bf 56} (2017) 56:18.

\bibitem{Zha.JDG.110-2018.543}
\leavevmode\vrule height 2pt depth -1.6pt width 23pt, \emph{Existence of solutions to the even dual Minkowski problem}, J. Differential Geom., {\bf 110} (2018) 543-572.

\bibitem{Zhu.Adv.262-2014.909}
G. Zhu, \emph{The logarithmic Minkowski problem for polytopes}, Adv.
  Math., {\bf 262} (2014) 909-931.

\bibitem{Zhu.JDG.101-2015.159}
\leavevmode\vrule height 2pt depth -1.6pt width 23pt, \emph{The centro-affine Minkowski problem for polytopes}, J. Differential Geom., {\bf 101} (2015) 159-174.


\end{thebibliography}




\end{document}
